% =====================================================================================
%  Big Data (ULPGC, 2026-2027) -- Individual Assignment %%NN%%
%  Author: Daniel Rodríguez Alonso
%  Build:  cd report && latexmk -pdf main.tex      (or: make report from the folder root)
%  Every number in this file must come from report/generated/*.tex or be traceable to results/.
% =====================================================================================
\documentclass[conference,a4paper]{IEEEtran}
\IEEEoverridecommandlockouts

\usepackage[utf8]{inputenc}
\usepackage[T1]{fontenc}
\usepackage{amsmath,amssymb}
\usepackage{graphicx}
\usepackage{booktabs}
\usepackage{siunitx}
\sisetup{group-separator={\,}, group-minimum-digits=4}
\usepackage[caption=false,font=footnotesize]{subfig}
\usepackage{xcolor}
\usepackage{listings}
\usepackage{microtype}
\usepackage{url}
\usepackage[hidelinks]{hyperref}
\usepackage[capitalise,noabbrev]{cleveref}

\lstset{basicstyle=\ttfamily\footnotesize, columns=fullflexible, breaklines=true,
        frame=single, numbers=none, keepspaces=true}
\graphicspath{{figures/}}

% Key numbers generated by the analysis scripts (\newcommand{\Something}{value})
\IfFileExists{generated/numbers.tex}{\input{generated/numbers.tex}}{}

% Draft helper: remove all \todo{} before submission (report-review checks this)
\newcommand{\todo}[1]{\textcolor{red}{[TODO: #1]}}

\begin{document}

\title{%%TITLE%%\\[0.3em]
\large Big Data -- Individual Assignment %%NN%%}

\author{\IEEEauthorblockN{Daniel Rodríguez Alonso}
\IEEEauthorblockA{Grado en Ciencia e Ingeniería de Datos\\
Universidad de Las Palmas de Gran Canaria (ULPGC)\\
Big Data, Academic Year 2026--2027}}

\maketitle

% ------------------------------------------------------------------ Abstract (150-200 words)
\begin{abstract}
\todo{Question. Method (implementations, inputs, sizes, machine, protocol). Main quantitative result
(with variability). Scope of the conclusion (this machine, these settings).}
\end{abstract}

\begin{IEEEkeywords}
matrix multiplication, benchmarking, performance evaluation, \todo{assignment-specific keywords}
\end{IEEEkeywords}

% ------------------------------------------------------------------ Introduction
\section{Introduction and Objective}
\label{sec:intro}
\todo{Motivation (why matrix multiplication matters for data-intensive computing).
The focused research question (RQ), the factor varied and the outcome measured.
Hypotheses, explicitly labelled as predictions (H1, H2, ...). Short list of contributions of this study.
If this assignment reuses earlier materials, say which (assignment, commit/tag).}

% ------------------------------------------------------------------ Background
\section{Background and Related Work}
\label{sec:background}
\todo{Algorithms/techniques needed to understand the study, each supported by a verified citation
(cite each with \textbackslash cite\{key\}). Cost model (e.g. $2n^3$ floating-point operations for the classical algorithm) and the
system effects relevant to the RQ (memory hierarchy, runtime, parallelism, communication).}

% ------------------------------------------------------------------ Methodology
\section{Methodology}
\label{sec:method}

\subsection{Implementations}
\todo{Each method: language, algorithm, loop order, storage layout, numeric type, what is own code and what
is library code. Short listing or pseudo-code if it helps.}

\subsection{Inputs}
\todo{Generator, seed(s), value range, numeric type, shared input files and checksums; for sparse inputs:
density levels, distribution of nonzeros, representation.}

\subsection{Experimental environment}
\begin{table}[t]
\centering
\caption{Execution environment (captured by \texttt{capture\_env.sh}).}
\label{tab:env}
\small
\begin{tabular}{@{}ll@{}}
\toprule
Item & Value \\
\midrule
CPU & \todo{model, cores/threads, caches} \\
Memory & \todo{RAM} \\
OS & \todo{distribution, kernel} \\
Compilers / runtimes & \todo{gcc + flags, JDK + JVM flags, Python} \\
Libraries & \todo{versions} \\
Thread settings & \todo{OMP\_NUM\_THREADS, BLAS threads, affinity} \\
\bottomrule
\end{tabular}
\end{table}

\subsection{Correctness validation}
\todo{Identity, zero and known small cases; comparison against a trusted reference; tolerances
(absolute and relative) and why they are appropriate; result of the tests.}

\subsection{Measurement protocol and budgets}
\todo{Timer and timing boundary (what is inside/outside), warm-ups (number, why), measured repetitions $k$,
statistic (median, IQR), memory metric and tool, declared size sequence, runtime and memory budgets,
stopping rule, interleaving/order of runs, machine-wide lock, exclusion rules.}

% ------------------------------------------------------------------ Results
\section{Results}
\label{sec:results}
\todo{Only observations here, no explanations. Every figure/table referenced in the text, with units,
$k$ and variability. Include failures/timeouts as results.}

% \begin{figure}[t]
%   \centering
%   \includegraphics[width=\columnwidth]{runtime_vs_n}
%   \caption{Median runtime vs.\ matrix size $n$ (log--log). Error bars: IQR over $k=5$ measured
%   repetitions. Kernel timing boundary. Machine: \cref{tab:env}.}
%   \label{fig:runtime}
% \end{figure}
% \IfFileExists{generated/table_summary.tex}{\input{generated/table_summary.tex}}{}

% ------------------------------------------------------------------ Discussion
\section{Discussion}
\label{sec:discussion}
\todo{Answer the RQ explicitly. Check each hypothesis against the evidence (supported / not supported).
Plausible explanations supported by evidence (and say which are conjectures). Threats to validity and
limitations: single machine, laptop frequency scaling, measurement tool limitations, sizes not tested.}

% ------------------------------------------------------------------ Conclusions
\section{Conclusions}
\label{sec:conclusions}
\todo{What this study establishes within its tested scope; what remains unresolved; future work
(link to the next assignment where relevant).}

% ------------------------------------------------------------------ Materials + AI statement
\section*{Reproducibility Materials}
\todo{Repository URL (and commit/tag of the submitted version), contents of the materials archive,
one-line command to reproduce the experiments and the figures, and where the raw data are.}

\section*{Use of Generative AI}
\todo{Course policy: disclose the tool, the purpose and the affected parts. Example: ``Claude Code (Anthropic,
model \ldots) was used to assist with \ldots in \ldots. All code was reviewed and executed by the author; all
measurements were produced by the author's runs and all references were verified by the author.'' Keep it
accurate and specific.}

% ------------------------------------------------------------------ References
\bibliographystyle{IEEEtran}
\bibliography{references}

\end{document}
